\begin{proof}[Proof of \cref{obj:lem:dorn-rate-scope}]
\begin{enumerate}
\item Fix a finite \(\gamma_0>1\), the source family \(\mathfrak P\), its common constants, and its selected propensity representatives. The design, moment, variance, Gaussian, smoothness, and global-tail hypotheses give Dorn's Assumption 2, including the restrictions retained from Assumption 1. The family-level condition in \cref{obj:def:dorn-a3} gives Assumption 3 on the original cube.

Dorn's fixed-setup upper comparison \citep[Section 2, Assumptions 1--3, Theorem 1(ii)]{Dorn2026GlobalOverlap} therefore supplies an estimator \(\widehat\mu_n\), measurable as a function of the sample for every \(n\) and \(x\), with comparison scale
\[
r^*_{\mu,n}
=
n^{-\beta/(2\beta+d\gamma_0/(\gamma_0-1))}
\qquad(n\ge1),
\qquad r^*_{\mu,0}=0.
\]
Its admissibility also gives measurability, for every \(n\), \(P\in\mathfrak P\), and \(R\in\mathbb R\), of
\[
\left\{
\|\widehat\mu_n-\mu_P\|_{L^\infty(P)}
>
\max\{Rr^*_{\mu,n},0\}
\right\}.
\]
For every \(\varepsilon>0\), the cited comparison supplies a finite \(c(\varepsilon)>0\) such that
\[
\limsup_{n\to\infty}\sup_{P\in\mathfrak P}
P^{\otimes n}\!\left\{
\|\widehat\mu_n-\mu_P\|_{L^\infty(P)}
>
c(\varepsilon)r^*_{\mu,n}
\right\}
\le\varepsilon.
\]
The setup was fixed before selecting the estimator and threshold constant, so both may depend on that setup. This proves the upper comparison.
% lean: CausalSmith.Stat.WeakOverlap.DornSourceRateTranscription

\item Fix \(\gamma_0>1\). The hypotheses \(d\ge1\), \(\beta>0\), and \(L_0>0\) permit application of \cref{obj:lem:arbitrary-family-lower-obstruction}. Its known-regression conclusion gives, for every nonempty family with selected treated curves equal pointwise to \(\mu^\circ\), every evaluation point \(x_0\), every sample size \(n\), and every \(c,r>0\),
\[
\inf_{T\ \mathrm{sample\text{-}measurable}}
\sup_{P\in\mathfrak P}
P^{\otimes n}\!\left\{
|T(\mathcal D_n)-\mu_P(x_0)|>cr
\right\}
=0,
\]
with the minimax value taken in the extended nonnegative reals.

The singleton conclusion of the same result supplies \(q'>3\) and \(M',C'>0\) for the independent experiment displayed in the statement. It supplies the uniform design, selected propensity \(1/2\), moment and regression-variance bounds, zero treated and control curves with the required source smoothness, global-tail bound, Gaussian conditional outcomes, and the family-level condition of \cref{obj:def:dorn-a3}. The conditional Gaussian distributions have variance one, hence
\[
\operatorname{Var}_{P}(Y\mid X,A)=1=\sigma^2
\qquad P\text{-almost surely},
\qquad \sigma=1.
\]
Thus the conditional-variance restriction also holds for the asserted parameter tuple. Finally, the singleton conclusion at \(x_0=0\) gives
\[
\inf_{T\ \mathrm{sample\text{-}measurable}}
P^{\otimes n}\!\left\{|T(\mathcal D_n)|>cr\right\}
=0
\qquad(n\in\mathbb N,\ c,r>0).
\]
These are the two pointwise conclusions.
% lean: CausalSmith.Stat.WeakOverlap.arbitraryFamily_lower_obstruction
\end{enumerate}
\end{proof}
